\begin{abstract}

\noindent
The power delivery network (PDN) of a multilayer printed circuit board is a
primary determinant of a product's electromagnetic compliance: cavity
resonances between power and reference planes raise the impedance seen by the
integrated circuits, converting switching noise into radiated emission and
jeopardizing approval in regulatory testing. Assessing this impedance currently
relies on full-wave electromagnetic solvers whose computational cost --- tens of
minutes per configuration --- precludes systematic exploration of the design
space during early layout, precisely when low-cost design decisions are still
available. This proposal develops and validates a physics-informed surrogate
model that predicts the complete PDN impedance curve $Z(f)$ from 1\,MHz to
1\,GHz, given eight geometric and material stackup parameters, with inference
latency below one millisecond. The work uses the six-layer PDN subset of the
SI/PI-Database maintained by Hamburg University of Technology, comprising 985
Latin-hypercube-sampled configurations simulated with a physics-based method
over 334 frequency points. The central contribution is the incorporation, into
the neural network loss function, of electromagnetic constraints available in
closed form --- the capacitive behavior governing the quasi-static regime, the
monotonicity of the response up to the series null, and the passivity condition
of the system --- in order to compensate for sample scarcity, a regime in which
purely data-driven models tend to violate the physics of the problem. A
preliminary investigation already carried out on the dataset confirmed the
validity of the capacitive anchor, whose measured slope is
$-1.018 \pm 0.008$ against a theoretical $-1$, and ruled out anchoring on modal
frequencies, which are not observable in the self-impedance within the analyzed
band. The model will be benchmarked against a suite of baselines (regularized linear
regression, random forests, gradient boosting, Gaussian processes and multilayer
perceptrons), embedded in a multi-objective evolutionary optimizer to search for
stackups with reduced electromagnetic interference risk, and interpreted through
symbolic regression and Shapley values, with the aim of distilling
engineer-readable design guidelines. Expected deliverables include a model
achieving a coefficient of determination above 0.90 for maximum impedance
prediction on a held-out test set, a set of interpretable analytical expressions
relating stackup parameters to electromagnetic risk, and a reproducible
open-source implementation documented according to spec-driven development
practices.

\end{abstract}
